%%%%%%%%%%%%%%%%%%%%%%%%%%%%%%%%%%%%%%%%%%%%%%%%%%%%%%%%%%%%%%%%%%%%%%%%%%%%%%%
%% Apêndice A -- Normas de funcionamento do laboratório
%%%%%%%%%%%%%%%%%%%%%%%%%%%%%%%%%%%%%%%%%%%%%%%%%%%%%%%%%%%%%%%%%%%%%%%%%%%%%%%
\apendice[Normas de Funcionamento do Laboratório]%
        {Normas de Funcionamento do Laboratório\\ de Sistemas Digitais e Microprocessadores}

A supervisão dos laboratórios, visando à aprendizagem e por questões de segurança,
estabelece para as aulas práticas e demais atividades no Laboratório de Eletrônica
Analógica e Digital um conjunto de Normas para Aulas Práticas e de Segurança, listadas
a seguir.

\secaoguia{I -- Normas gerais}

\begin{normas}[I]
  \norma{OBJETOS} --- os objetos volumosos de uso pessoal (mochilas, pastas, livros,
  cadernos etc.) devem ser colocados nos escaninhos disponíveis, na frente da sala ou
  sob as bancadas do laboratório, mas nunca sobre as mesmas;

  \norma{CELULAR} --- NÃO é permitido o uso do telefone celular ligado durante as aulas
  práticas. Recomenda-se desligá-lo antes de cada aula;

  \norma{PONTUALIDADE} --- para as aulas de laboratório existe uma tolerância de até 15
  minutos. \textbf{O aluno que chegar após este período não poderá assistir às aulas e
  só poderá efetuar a sua reposição por motivo justificável (por escrito).}

  \norma{NÚMERO DE ALUNOS POR BANCADA} --- as montagens dos circuitos serão feitas por
  até 4 (quatro) alunos por bancada, divididos em grupos de dois alunos, que deverão
  fazer o relatório da aula por grupo e entregá-lo na semana seguinte à aula;

  \norma{POSTURA PROFISSIONAL:}
  \begin{enumerate}[label=\alph*),leftmargin=1.2cm]
    \item não será permitido aos alunos assistirem aula trajando bermuda e/ou
          camisetas, por medidas de segurança e por aspectos de postura profissional;
    \item o laboratório não é um lugar de bate-papo ou de perda de tempo;
    \item não é permitido ao aluno fumar ou fazer lanche no laboratório;
  \end{enumerate}
\end{normas}

\secaoguia{II -- Normas de organização e segurança}

\begin{normas}[II]
  \norma{PROCEDIMENTO INICIAL} --- antes de cada montagem, ligue o disjuntor
  diferencial da bancada. Após o término da mesma, desligue-o;

  \norma{DISPOSIÇÃO DOS COMPONENTES E EQUIPAMENTOS} --- ao término de cada montagem, os
  equipamentos e componentes devem ser dispostos como recebidos antes do início da aula;

  \norma{PARA QUALQUER PROCEDIMENTO} --- execute-o somente se estiver habilitado. Na
  DÚVIDA, informe-se com o professor ou com o laboratorista;

  \norma{SOLICITE A PRESENÇA DO PROFESSOR} em sua bancada, no caso de dúvidas ou de
  acidentes.
\end{normas}
